\subsection{Flujo Máximo / Corte Mínimo (Dinic) para Etiquetado Binario}

\graybold{Preguntas guía}

\begin{itemize}
\item ¿Necesito asignar a cada elemento una de dos etiquetas (0/1, s/t, bajo/alto) minimizando un costo que combina un costo individual por elemento y un costo por cada par de elementos vecinos con etiquetas distintas?
\item ¿El costo por vecinos "diferentes" es simétrico y no depende de cuál de los dos valores tomó cada uno, solo de si difieren?
\item ¿El problema se puede reformular como "separar los nodos de un grafo en dos grupos" minimizando el peso total de las conexiones que cruzan la separación (corte)?
\item ¿La función de costo es submodular (no hay ninguna ventaja "extra" por combinaciones específicas de etiquetas más allá del costo por diferir)?
\end{itemize}

\graybold{¿Para qué sirve?} Muchos problemas de decisión binaria con costos de "cambio" individuales y costos de "frontera" entre vecinos (segmentación de imágenes, problemas de selección de proyectos, particiones de grafos con penalización por aristas cortadas) se pueden resolver de forma exacta modelando un grafo con una fuente $S$ y un sumidero $T$, donde cada elemento es un nodo, y el corte mínimo entre $S$ y $T$ corresponde exactamente al costo mínimo de la asignación óptima. El algoritmo de Dinic calcula flujo máximo (que por el teorema de flujo máximo–corte mínimo es igual al corte mínimo) en \bigO{V^2 E} en el peor caso, aunque en la práctica —y en particular en grafos tipo grilla como este— es mucho más rápido.

\graybold{Complejidad:} \bigO{V^2 E} en el peor caso general para Dinic, donde $V$ es el número de nodos y $E$ el número de aristas del grafo de flujo construido.

\graybold{Fundamento teórico:} Se construye un grafo con un nodo por celda, más una fuente $S$ (representa la etiqueta "bajo") y un sumidero $T$ (representa "alto"). Para cada celda $v$:
\begin{itemize}
\item Si originalmente es baja (\texttt{.}), se agrega la arista $S \to v$ con capacidad $B$: esta arista solo forma parte del corte si $v$ termina del lado de $T$ (es decir, cambió a alta), cobrando exactamente el costo de cambio.
\item Si originalmente es alta (\texttt{\#}), se agrega la arista $v \to T$ con capacidad $B$: solo forma parte del corte si $v$ queda del lado de $S$ (cambió a baja).
\end{itemize}
Para cada par de celdas vecinas $u, v$ (horizontal o verticalmente adyacentes), se agregan las dos aristas dirigidas $u \to v$ y $v \to u$, ambas con capacidad $A$. Como el corte solo cuenta aristas que van del lado de $S$ al lado de $T$, exactamente una de las dos direcciones contribuye al corte cuando $u$ y $v$ terminan en lados distintos, cobrando $A$ sin importar cuál de las dos quedó en cada lado —esto es lo que hace que la construcción capture correctamente un costo \emph{simétrico} por vecinos diferentes.

Por el teorema de flujo máximo–corte mínimo, el flujo máximo entre $S$ y $T$ es igual al peso del corte mínimo, que por construcción es exactamente el costo mínimo total (suma de costos de cambio más costos de frontera) sobre todas las asignaciones binarias posibles. Esta equivalencia es válida porque la función de costo es submodular: no existe combinación de etiquetas vecinas que sea "más barata de lo esperado" fuera de lo que la construcción ya captura, condición necesaria para que la relajación por corte mínimo coincida exactamente con el óptimo combinatorio (a diferencia de un problema con más de dos etiquetas o costos no submodulares, donde el corte mínimo podría no ser exacto).

\graybold{Ejercicio aplicado}

Dado un campo de $N \times M$ parcelas, cada una originalmente en un nivel bajo (\texttt{.}) o alto (\texttt{\#}), con costo $A$ por cada par de parcelas vecinas (4-adyacentes) que terminen en niveles distintos y costo $B$ por cambiar el nivel de una parcela, hallar el costo mínimo total para que el tractor recorra el campo.

\graybold{I/O:} Se usó lectura total con tokenización (\texttt{sys.stdin.buffer.read().split()}) combinada con el patrón de "hasta agotar los datos" (sin contador de casos): el archivo trae varios casos de prueba consecutivos sin un $T$ inicial ni líneas en blanco que preservar, y cada fila del campo es un token sin espacios internos, así que tokenizar todo de una vez y consumir con un índice cubre tanto los enteros como las filas del grid de forma uniforme. El tamaño por caso es pequeño ($N,M \le 50$, hasta 2500 celdas), así que se probó el rendimiento con varios casos consecutivos de $50\times 50$: un único caso tardó 0.22s, y 100 casos consecutivos de $50\times 50$ tardaron 1.47s en total, tiempos holgados dado que el enunciado no fija una cota explícita para el número de casos por archivo.

\begin{lstlisting}[language=Python]
import sys
from collections import deque


class Dinic:
    def __init__(self, n):
        self.n = n
        self.graph = [[] for _ in range(n)]  # cada arista: [destino, capacidad, indice_reversa]

    def add_edge(self, u, v, cap):
        self.graph[u].append([v, cap, len(self.graph[v])])
        self.graph[v].append([u, 0, len(self.graph[u]) - 1])

    def bfs(self, s, t):
        self.level = [-1] * self.n
        self.level[s] = 0
        q = deque([s])
        while q:
            u = q.popleft()
            for e in self.graph[u]:
                v, cap, _ = e
                if cap > 0 and self.level[v] < 0:
                    self.level[v] = self.level[u] + 1
                    q.append(v)
        return self.level[t] >= 0

    def _dfs_una_ruta(self, s, t):
        # DFS iterativo (evita limite de recursion) que busca UNA ruta de
        # aumento en el grafo de niveles actual, usando punteros it[] que
        # persisten entre llamadas dentro de la misma fase para saltar
        # rapido las aristas ya agotadas o que llevan a callejones sin salida
        stack = [s]
        camino = []  # lista de (nodo, indice_arista) recorridos
        while stack:
            u = stack[-1]
            if u == t:
                return camino
            avanzo = False
            while self.it[u] < len(self.graph[u]):
                v, cap, _ = self.graph[u][self.it[u]]
                if cap > 0 and self.level[v] == self.level[u] + 1:
                    camino.append((u, self.it[u]))
                    stack.append(v)
                    avanzo = True
                    break
                self.it[u] += 1
            if not avanzo:
                self.level[u] = -1  # callejon sin salida, no volver a visitarlo
                stack.pop()
                if camino:
                    camino.pop()
        return None  # no quedan rutas de aumento en este grafo de niveles

    def dfs_blocking_flow(self, s, t):
        self.it = [0] * self.n
        total = 0
        while True:
            camino = self._dfs_una_ruta(s, t)
            if camino is None:
                break
            f = min(self.graph[u][i][1] for u, i in camino)
            for u, i in camino:
                e = self.graph[u][i]
                e[1] -= f
                rev = self.graph[e[0]][e[2]]
                rev[1] += f
            total += f
        return total

    def max_flow(self, s, t):
        flow = 0
        while self.bfs(s, t):
            flow += self.dfs_blocking_flow(s, t)
        return flow


def resolver(n, m, a, b, grid):
    S = n * m
    T = n * m + 1
    din = Dinic(n * m + 2)

    def nodo(r, c):
        return r * m + c

    for r in range(n):
        for c in range(m):
            u = nodo(r, c)
            if grid[r][c] == '.':
                # cambiar esta celda a alta cuesta B: se paga si termina
                # del lado de T (separada de la fuente)
                din.add_edge(S, u, b)
            else:
                # cambiar esta celda a baja cuesta B: se paga si termina
                # del lado de S
                din.add_edge(u, T, b)
            # aristas con la celda de la derecha y de abajo (evita duplicar)
            if c + 1 < m:
                v = nodo(r, c + 1)
                din.add_edge(u, v, a)
                din.add_edge(v, u, a)
            if r + 1 < n:
                v = nodo(r + 1, c)
                din.add_edge(u, v, a)
                din.add_edge(v, u, a)

    return din.max_flow(S, T)


def main():
    data = sys.stdin.buffer.read().split()
    idx = 0
    salidas = []
    while idx < len(data):
        n = int(data[idx]); m = int(data[idx + 1])
        a = int(data[idx + 2]); b = int(data[idx + 3])
        idx += 4
        grid = []
        for _ in range(n):
            grid.append(data[idx].decode())
            idx += 1
        salidas.append(str(resolver(n, m, a, b, grid)))
    sys.stdout.write('\n'.join(salidas) + '\n')


if __name__ == "__main__":
    main()
\end{lstlisting}
